\section{Beats: Time Without a Clock}

The rule updates a vibe ``at each moment.'' What is a moment, here? This chapter is about how time works in the theory at the smallest scale: the beat, the discrete tick of update, and why there is no single universal clock ticking for the whole universe.

\subsection{Time comes in ticks}

Time in this theory is not a smooth flowing river. It comes in discrete steps, called beats. A beat is one round of updating: vibes apply the rule, tones change, and that is one tick. Then another beat, and another. Between beats nothing is happening. The world is a sequence of still frames, and time is the flicking from one frame to the next.

This matches the discreteness everywhere else in the theory. The states are discrete (three tones). The space is discrete (a graph of cells). And time is discrete too (beats). There are no infinitely-fine moments, just as there are no infinitely-fine positions. The universe is grainy all the way down, in space and in time.

\subsection{No global clock}

Here is the surprising part. There is no single master clock that ticks for the whole universe at once. Vibes update locally, each in its own neighborhood, without waiting for a universe-wide signal that says ``now.'' Updates happen here and there, in their own order, with no global ``this instant'' binding them all together.

Why forbid a universal now? Because a universal now would be a preferred frame, and that is exactly what Einstein's relativity rules out and experiment confirms. If there were one true ``now'' slicing the whole universe, you could in principle detect it, and physics says you cannot. Different observers moving differently disagree about which distant events happen ``at the same time,'' and none of them is right or wrong. There is no master slicing. So the theory refuses to build one in. Updating is local, and a global ``now'' is not part of the machinery.

\subsection{Then what is time?}

If there is no master clock, what makes time real and gives it a direction? Two things, both already in our hands.

Time is the depth of before-and-after. Even without a master clock, there is a clear order of cause: this vibe's state was set using that neighbor's earlier state, which was set from its neighbors' earlier states, and so on back. Follow those chains of ``this came from that'' and you get a web of before-and-after relationships. The length of the longest such chain up to some event is, in effect, how much time has passed to reach it. Time is the depth of the causal web, not the reading of a clock.

And time has a direction because the crystal only grows. New cells are added at the frontier, never removed. Relationships accumulate and are never undone. The past, the deep interior of the crystal, is fixed and finished. The future, the growing edge, is open and being made. The arrow of time is simply that the heap of what-has-happened only ever gets bigger. We saw this back in the bootstrap: distinctions can be added but not subtracted, and that one-way pile-up is time itself.

\subsection{The present is the edge}

Put it together and you get a vivid picture. The present moment is the growing frontier of the crystal, the freshly added cells where the action is. The past is everything already grown, frozen behind. The future does not exist yet. It is being made, beat by beat, at the edge. There is no master clock because the ``now'' is not a slice through a finished block, it is the living rim of a crystal that is still growing. You are not moving through time. You are part of the edge where time is being made.

\textbf{Takeaway:} time comes in discrete beats, ticks of updating, with no universal clock, because a single ``now'' would be a forbidden preferred frame. Time is really the depth of the before-and-after web, and its arrow is the one-way growth of the crystal. The present is the growing edge.
